\section{Linear-group growth and the affine Lehmer model}
\label{sec:growth}

The growth strategy seeks a gap at fixed matrix dimension, uniform over the subgroup and finite generating set.
The affine model gives a concrete return to the scalar objective.
Its subsidiary strategy studies all original orbit collisions and their complete weights.

\begin{definition}\label{gr:word-growth}
For a nonempty finite $S\subset GL_r(\C)$, let $S^n$ contain products of exactly $n$ elements and set
\[
 \rho(S)=\lim_{n\to\infty}\frac1n\log|S^n|.
\]
\end{definition}
\begin{proposition}\label{gr:gap-equivalence}
The inequality $\mathfrak L>0$ is equivalent to the following statement:
for each fixed $r\ge1$ there is $\delta_r>0$ such that every nonempty finite $S\subset GL_r(\C)$ satisfies $\rho(S)=0$ or $\rho(S)\ge\delta_r$.
\end{proposition}
This equivalence uses the structural reduction of \citet[Theorem~4]{BV20}.
For $\alpha\ne0$, the particular family
\[
 S_\alpha=\left\{
 \begin{pmatrix}\alpha&1\\0&1\end{pmatrix},
 \begin{pmatrix}\alpha&-1\\0&1\end{pmatrix}\right\},
\]
already has matrix dimension two, although the algebraic degree of $\alpha$ is unrestricted.
A sequence in increasing matrix dimensions would not negate a fixed-dimension gap.

\begin{proposition}\label{gr:entropy-comparison}
For a nonzero algebraic integer $\alpha$, the original binary rates of Definition~\ref{fo:binary-law} satisfy the following comparison, with an absolute $c>0$:
\begin{equation}\label{fo:entropy-order}
 c\min\{m,\log2\}\le h_1\le h_0=\rho(S_\alpha)\le m,
 \qquad h_2\le h_1.
\end{equation}
\end{proposition}
The lower bound uses \cite[Theorem~5]{BV20}; signs and bits give the same fibers by an invertible affine transformation.
The last bound also has a direct full-embedding norm-grid proof: the number of possible sums is at most $(4n)^d\exp((n-1)m)$.


\subsection{Original-orbit collisions and affine binary growth}
\label{sec:returns}

The binary affine model has a useful exact difference process.
Let $\alpha$ be an algebraic unit, $\mathcal O=\Z[\alpha]$, and let $\eta_j=\varepsilon_j-\varepsilon'_j$, where the two binary families are independent and fair.
Thus
\[
 \Prob(\eta_j=0)=\tfrac12,\qquad
 \Prob(\eta_j=1)=\Prob(\eta_j=-1)=\tfrac14.
\]
The probability $p_n=\Prob(\sum_{j<n}\eta_j\alpha^j=0)$ is the collision probability of the binary sum.
Equivalently it is the zero-return probability of $X_{j+1}=\alpha X_j+\eta_j$, started at zero.
Reversing the independent digit order identifies these two conventions.
All probabilities in this subsection use this law unless a change is expressly specified.

\subsubsection{An infinite core for the original Lehmer number}
\label{ret:core}

A Salem number is a real algebraic unit $\beta>1$ whose other conjugates are $\beta^{-1}$ and points of the unit circle, with at least one such unit-circle conjugate \cite[\S8]{Smy08}.
\begin{definition}\label{ret:return-core}
The \emph{return core} $\mathcal D\subseteq\mathcal O$ consists of states lying on finite paths from zero back to zero.
It records possible paths, whereas $p_n$ counts them with their original weights.
This distinction remains important even when the core has an exact geometric description.
\end{definition}

\begin{theorem}\label{ret:lehmer-core}
Let $\beta>1$ be the expanding root of
\[
 P_L(t)=t^{10}+t^9-t^7-t^6-t^5-t^4-t^3+t+1.
\]
Put $B=(\beta-1)^{-1}$, $S=\beta/(\beta-1)$, and write $x_+=x(\beta)$, $x_-=x(\beta^{-1})$ for $x\in\Z[\beta]$.
Then the original return core is
\begin{equation}\label{ret:strip}
 \mathcal D=\{x\in\Z[\beta]: |x_+|<B,\ |x_-|<S\}.
\end{equation}
If $\zeta_1,\ldots,\zeta_4$ are one representative from each conjugate pair on the unit circle and $\rho(x)^2=\sum_{j=1}^4|x(\zeta_j)|^2$, then
\[
 cR^8\leq\#\{x\in\mathcal D:\rho(x)\leq R\}
       \leq C(R+1)^8\qquad(R\geq R_0)
\]
for positive constants depending on this fixed polynomial.
\end{theorem}

The main input to this theorem is a complete co-return certificate: every actual state with $|x_+|\leq1$ and $|x_-|\leq S$ reaches zero.
It was proved by a joint eight-dimensional neutral descent, followed by a complete finite residue check.
For a digit word $w$ of length $\ell$, its neutral displacement before the endpoint rotation is
\[
 c_w=\left(\sum_{i=1}^{\ell}\eta_i\zeta_j^{-i}\right)_{j=1}^4.
\]
The normalized vectors $c_w/\lVert c_w\rVert^2$, rather than the unnormalized displacements, give the useful descent certificate.
On each of the expanding intervals $[1/2,3/4]$ and $[3/4,1]$, exact convex containment certificates using words of length at most eleven imply a squared-radius decrease of at least $57/100$ when $\rho\geq8$.
The complete compact enumeration contains $797507$ candidate states; exact paths and certified exclusions cover it.
Separate computations checked the containment, complete enumeration, paths and final assembly.
These are certificates for $P_L$ itself.
Mercat's earlier descent and finite-residue method concerns the different quartic $t^4-2t^3+t^2-2t+1$ \cite[Proposition~4.23]{Mer13}; its numerical certificate is not being transferred to $P_L$.

The passage from narrow co-return to \eqref{ret:strip} is elementary.
While $1<|x_+|<B$, choose the digit opposite to its sign; then $|x_+|$ changes to $\beta|x_+|-1$ and reaches the narrow strip in finitely many steps.
The stable bound is preserved.
Reciprocal conjugation, followed by multiplication by $\beta^{-1}$, reverses zero-ending words and supplies reachability from zero.
Conversely, the geometric sums along a finite incoming and outgoing word give the two strict inequalities in \eqref{ret:strip}.
The state count follows from the full embedding lattice, using a cylinder for the upper bound and differences of coefficient-box points for the lower bound.
It is not a count of equally probable words.

A later control argument improves the original selected-word estimate to linear length.
\begin{proposition}\label{ret:linear-control}
For the fixed $\beta$ of Theorem~\ref{ret:lehmer-core}, every state with $|x_+|\le1$ and $|x_-|\le S$ has a zero-ending word of length at most $A_\beta\rho(x)+B_\beta$, for finite constants depending only on $\beta$.
\end{proposition}
A finite directional cover with nonzero endpoint margins gives a fixed radial decrease per bounded-length block; the initial zero prefix costs only $O_\beta(\log(1+\rho))$ by the integer norm.
The finite termination certificate above closes the argument.
The same control is conditional on compact termination for a general Salem number below the golden ratio.

This connects selected paths with an all-path statement.
For a fixed Salem number $1<\beta<2$ of degree $2s+2$, define $\rho$ using its $s$ unit-circle conjugate pairs and set
\[
 \gamma_\beta=\frac{\log2\log\beta}{\log(1+\beta)},\qquad
 H(x)=\Prob_x(\tau<\infty),\quad
 \tau=\inf\{n\ge1:X_n=0\}\quad(x\ne0).
\]
\begin{proposition}\label{ret:hitting-rate}
There is a finite $D_\beta>0$ such that every nonzero $x\in\Z[\beta]$ satisfies
\begin{equation}\label{ret:hitting}
 H(x)\le\min\{1,D_\beta
           e^{-\gamma_\beta\lceil\rho(x)/\sqrt{s}\rceil}\}.
\end{equation}
If every state in the narrow strip $|x(\beta)|\le1$, $|x(\beta^{-1})|\le\beta/(\beta-1)$ admits a zero-ending word of length at most $A_\beta\rho(x)+B_\beta$, then on that strip $H(x)\ge e^{-a_\beta\rho(x)-b_\beta}$ and $\E_x[\tau\mid\tau<\infty]=O_\beta(1+\rho(x))$.
\end{proposition}
\begin{proof}[Proof outline]
The maximum-atom collision estimate and the necessary length $n\ge\rho(x)/\sqrt{s}$ give the upper bound.
An admissible selected word of linear length has original probability at least its $4^{-\text{length}}$ weight, giving the lower bound.
Dividing the complete hitting tail by that lower bound gives a full exponential conditional tail after an input-dependent linear time and hence the conditional mean estimate.
\end{proof}
This conditional law is different from the original unconditioned walk.
The constants, and already $\gamma_\beta$, do not yield a positive uniform bound as $\beta\downarrow1$.

\subsubsection{A necessary diffusion scale for stationary laws}
\label{ret:small-ball}

Let $\alpha$ now be any algebraic unit with $s\ge1$ selected unit-circle embeddings, and put $Z(x)=(\sigma_1(x),\ldots,\sigma_s(x))\in\C^s$.
Consider an irreducible positive recurrent Markov chain on a subset of the actual return core containing zero, using only the original edges $x\mapsto\alpha x+\eta$.
If its transition probabilities are $q_x(\eta)$ and invariant probability is $\pi$, its information cost per step relative to the original law is
\[
 I=\sum_x\pi(x)\sum_{\eta=-1}^1q_x(\eta)
           \log\frac{q_x(\eta)}{\nu_\eta},
 \qquad(\nu_{-1},\nu_0,\nu_1)=(\tfrac14,\tfrac12,\tfrac14).
\]
Zero terms have value zero.
The comparison below includes the full original paths, even when they leave the recurrent subset.

\begin{proposition}\label{ret:stationary-radius}
For every such chain, every $c\in\C^s$ and $r\ge0$,
\begin{equation}\label{ret:smallball-bound}
 \pi\{\lVert Z(x)-c\rVert\le r\}
 \le 6\bigl[I(1+r^2/s)\bigr]^{1/4}+\sqrt{I/2}.
\end{equation}
No finite-support or moment assumption is needed.
In particular, a ball of mass at least $u>0$, with $0<I\le u^2/8$, must satisfy $r^2/s\ge u^4/(4096I)-1$.
\end{proposition}

The proof uses a dimension-free separated-sum argument of Greene and Kleitman \cite[Section~9]{GK78}.
For deterministic Hilbert-space vectors of common norm $a$ and the original independent digits,
\[
 \Prob\left\{\sum_{j=1}^n\eta_jv_j\in\overline B(c,r)\right\}
 \le\min\left\{1,\sqrt{\frac{24+96r^2/a^2}{n}}\right\}.
\]
Grouping into independent blocks gives a positive probability of a displacement longer than $r$; the separated-sum bound then controls their independently chosen signs.
In the neutral embedding $a^2=s$.
Compare the stationary path of length $n$ with the original independent-digit process started from the same initial distribution $\pi$.
Their relative entropy is exactly $nI$.
Endpoint comparison adds $\sqrt{nI/2}$ to the displayed concentration estimate.
Optimizing the integer $n$ proves \eqref{ret:smallball-bound}.
If $I=0$, letting $n$ tend to infinity before optimization would give zero mass to every finite-radius ball, which is impossible for a probability law on these finite-dimensional neutral coordinates.

Thus, along any hypothetical $m(\alpha)\to0$ sequence with admissible laws satisfying $I\le2m(\alpha)$, every radius $r=o(\sqrt{s/m(\alpha)})$ has uniformly vanishing mass, uniformly in its center.
This is a necessary spatial scale, not a contradiction: an explicit stationary law for the torsion control $\alpha=-1$ attains the corresponding low-information scale.
A further arithmetic or original-return restriction is needed to exclude it for the relevant nontorsion inputs.

\subsubsection{What lower covariance certificates can establish}
\label{ret:blocks}

Fix a nontorsion algebraic unit $\alpha$ and a length $2n$ word.
For $A\subseteq\Omega=\{0,\ldots,2n-1\}$ define
\[
 U_A=\sum_{j\in A}\eta_j\alpha^j,\qquad
 Y_A=-\alpha^{-n}\sum_{j\in A,\ j<n}\eta_j\alpha^j,
 \qquad p_A=\Prob(U_A=0).
\]
On $U_A=0$, the second expression for $Y_A$ is $\sum_{j\in A,\ j\ge n}\eta_j\alpha^{j-n}$.
For $A=\Omega$ this is the original midpoint under the complete return condition.
Let $\Phi$ denote all real Minkowski coordinates, using real and imaginary parts once for each complex pair.

\begin{proposition}\label{ret:block-lower}
For pairwise disjoint nonempty supports $A$, conditional covariance gives the positive-semidefinite inequality
\begin{equation}\label{ret:block-certificate}
 \Cov(\Phi(Y_\Omega)\mid U_\Omega=0)
 \succeq\sum_Ap_A\Cov(\Phi(Y_A)\mid U_A=0).
\end{equation}
\end{proposition}
Indeed, condition first on all block endpoints.
Internal words then remain independent.
Keeping only each block's zero endpoint pays at least its original mass $p_A$, since a difference of two identical binary laws has its largest atom at zero.
Overlapping supports cannot be added without an additional averaging factor.
For the primitive quartic $t^4-t^3-t^2-t+1$, identify its neutral embedding with its two real coordinates $Z$.
This mechanism gives the actual neutral midpoint covariance bound
\[
 \Cov(Z(Y_\Omega)\mid U_\Omega=0)
       \succeq\frac{5-\sqrt{13}}{4096}I_2\qquad(n\ge4).
\]
This is a constant floor, not a diffusive lower bound.

\begin{proposition}\label{ret:bounded-certificate}
For every fixed nontorsion algebraic unit, the trace of the right side of \eqref{ret:block-certificate} is bounded independently of $n$ and of the disjoint supports.
Its neutral contribution from supports at distance at least $G$ from the midpoint cut decreases doubly exponentially in $G$.
The full contribution also has nonneutral exponential terms.
All constants depend on the algebraic input.
\end{proposition}

Here is the arithmetic reason for the bound.
Write $d_0$ for the real neutral dimension, $m=m(\alpha)$, and $r_i$ for all conjugate moduli, counting both complex partners in products.
Put $C=\prod_{r_i\ne1}|r_i-1|^{-1}$.
If a returned support has cardinality $\ell$, closest distance $g$ from the cut, and nonzero midpoint, the full integer norm gives
\begin{equation}\label{ret:block-norm}
 1\le|\Norm(Y_A)|\le Ce^{-2mg}(\ell/2)^{d_0}.
\end{equation}
Meanwhile, choosing an expanding modulus $r_*$ and $k=\min\{j\ge1:r_*^j>2\}$ gives $p_A\le e^{-a\ell}$, $a=(\log2)/k$, by injectivity of a separated binary subcube.
For $d_0>0$, \eqref{ret:block-norm} forces $\ell\ge2C^{-1/d_0}e^{2mg/d_0}$ whenever the midpoint contributes.
Its neutral squared norm is at most $d_0\ell^2/8$.
At most two disjoint supports have any prescribed closest distance, so summing these paid contributions proves the assertion.
When $d_0=0$ the neutral term is absent; geometric bounds at every expanding and contracting place directly bound the actual full covariance.

The exact improved coefficient in \eqref{ret:block-certificate} is $p_A\alpha_A$, where
\[
 \alpha_A=\frac{p_{\Omega\setminus A}}{p_\Omega},
 \qquad1\le\alpha_A\le2^{|A|}.
\]
If $\alpha_A\le Ke^{\gamma|A|}$ for all contributing blocks and some $\gamma<a$, the same argument still gives a bounded certificate.
Unbounded refined certificates therefore require positive neutral dimension and $\limsup\log\alpha_A/|A|\ge a$ along growing supports.
This necessary amplification is not sufficient: taking the whole word as one block already meets it, but merely reproduces the unknown actual covariance.
An upper bound on these lower certificates does not bound the actual neutral covariance from above.

\subsubsection{Localizing the Gaussian comparison loss}
\label{ret:neutral-rank}

An independent upper-comparison argument locates another precise remaining difficulty.
Let $X$ be an integral $d\times b$ matrix of full row rank, $Q=XX^T$, $S_X=X\eta$ and $p=\Prob(S_X=0)$.
For $0<W\preceq I_d$, set $a=\log2$, $D=(W^{-1}+Q)^{-1}$ and
\[
 G(u)=e^{au^TDu}\E e^{-a(S_X+u)^TW(S_X+u)}.
\]
\begin{proposition}\label{ret:gaussian-profile}
The proved original atom profile has coefficient $1/2$:
$\Prob(S_X=z)\le pe^{-z^TQ^{-1}z/2}$.
Completing the Gaussian leaves a positive mismatch:
\begin{align*}
 G(u)&\le K_d\det(W)^{-1/2}p\,e^{u^TR_0u},\\
 R_0&=a\bigl[(W^{-1}+Q)^{-1}-(W^{-1}+2aQ)^{-1}\bigr],
 \qquad K_d=2^{d/4}(2\pi/a)^{d/2}.
\end{align*}
One has $0\preceq R_0\preceq(a-1/2)Q^{-1}$.
\end{proposition}
\begin{proof}[Proof outline]
The original half-Gaussian atom profile is applied before completing the square.
Completing it against the factor with coefficient $a=\log2$ gives the stated resolvent difference.
Comparing the complete shifted lattice sum with its unit cubes gives $K_d\det(W)^{-1/2}$.
Conjugation by $Q^{1/2}$ gives the two positive-semidefinite comparisons for $R_0$.
\end{proof}

For the actual power-basis columns of a fixed nonzero algebraic integer, let $M$ be multiplication by that integer, let $v$ represent $1$, and put $A_i=[-M^dv,\ldots,-M^{d+i-1}v]$, with $A_0$ empty.
Use $X=[-M^{d+i}v,\ldots,-M^{d+i+b-1}v]$ and $W=(I+A_iA_i^T)^{-1}$, where $b\geq d$.
If $H_h$ is the full hyperbolic coordinate map, set
\[
 Q_h=H_hQH_h^T,\quad C_N=Q-QH_h^TQ_h^{-1}H_hQ,\quad
 R_N=a\bigl[D-(D^{-1}+(2a-1)C_N)^{-1}\bigr].
\]
Absent hyperbolic coordinates mean that the subtracted term is zero.
Let $d_0$ be the real neutral dimension.
The normalized projected energy of an endpoint $z$ is $(H_hz)^TQ_h^{-1}(H_hz)$, or zero when hyperbolic coordinates are absent.
For these actual columns, write $G_{i,b}$ for $G$ and note that $p=p_b$.

\begin{theorem}\label{ret:neutral-localization}
For every fixed nonzero algebraic integer, the entire hyperbolic projected support has normalized energy at most a finite $M_h$, independent of $i\ge0$ and $b\ge d$.
With $\delta=a-1/2$ one has
\[
 G(u)\le e^{\delta M_h}K_d\det(W)^{-1/2}p\,e^{u^TR_Nu},
 \qquad \operatorname{rank}R_N=d_0,\quad
 \ker R_N=D^{-1}\operatorname{range}(H_h^T).
\]
If $i_b=o(b)$ and $t_b\ge0$ with $t_b=o(b)$, every zero-containing set $\mathcal V_b\subset\{u:u^TR_Nu\le t_b\}$ satisfies
\[
 \lim_{b\to\infty}-\frac1b\log\sup_{u\in\mathcal V_b}G_{i_b,b}(u)=h_2.
\]
The complete prefix expense $\det(I+A_iA_i^T)^{1/2}$ is included.
After undoing prefix multiplication, pure-neutral translations of radius $o(b)$ also preserve this rate.
\end{theorem}
\begin{proof}[Proof outline]
On the full original support, pay the bounded projected energy once.
The resulting atom profile has coefficient $a$ on hyperbolic coordinates and $1/2$ on their complement.
Square completion gives $R_N$ exactly.
Whitening $Q$ identifies $C_N$ as a positive form of rank $d_0$; its resolvent difference has the same rank and the displayed kernel.
At fixed input, the full hyperbolic Gram matrices have a positive spectral floor after rescaling, uniformly in prefix and word length.
The prefix determinant has sublinear logarithm when $i=o(b)$, so the finite cylinder bound preserves the collision exponent.
The pure-neutral domain additionally uses the complete phase-Gram and cross-Schur bounds.
\end{proof}
Macroscopic neutral translations remain outside this conclusion.
The fixed-input limit does not identify $h_2$ with $m$ or supply input-uniform constants.


\paragraph{The remaining original-law input.}
The return-core theorems identify actual states and admissible words.
The probabilistic objective concerns all such words with their original weights.
An infinite core, a path to zero from each state, and a stationary law obtained by conditioning are different results, with different quantitative costs.
The current fixed-input estimates do not provide a positive entropy gap uniform in the algebraic input.

\begin{question}\label{out:neutral}
Which arithmetic constraints on the complete original return law force useful neutral-coordinate occupation or moment estimates, with constants that remain effective as total Mahler measure decreases and degree grows?
\end{question}
The returned-block analysis gives a concrete test for proposed lower bounds.
A disjoint certificate carrying only its original block mass has a fixed-input ceiling.
Conditioning may amplify that mass, but the exact complementary return ratio must be estimated.
Conditioning on the whole word already restores a tautological identity and supplies no growth estimate by itself.
On the Fourier side, a complete fractional return moment yields useful critical regularity, but that moment is an additional premise which still has to be proved for the relevant nontorsion process.


Here is the precise long-window limitation for one collective-column lower certificate.
Let $P$ be a monic irreducible minimal polynomial of a nonzero algebraic integer, with total measure $m>0$ and $r$ exterior conjugates, counting both members of a complex pair, and put $s=e^{-m}$.
Let $\mathcal B$ be the finite Blaschke product with zeros $1/\overline\alpha$ for these exterior conjugates, normalized by $\mathcal B(0)=s$.
Write $F=\mathcal B/s=\sum_{j\ge0}c_jz^j$ and its formal inverse $H=1/F=\sum_{j\ge0}h_jz^j$.
For an integer $b\ge1$, let $F_b=(c_{j-i})$ and $C_b=(h_{j-i})$ be upper triangular $b$-by-$b$ Toeplitz matrices, with coefficients of negative index zero, and put $D_b=C_b-I_b$.
For $A\ge\log2$, define
\[
 \begin{split}
 \kappa(A)&=\min\{\sqrt{\pi/A}\,e^{-\pi^2/\log2},
                         e^{-25A/16}(1-2^{-3/2})\},\\
 \gamma(A)&=\kappa(A)/32,\qquad K(A)=2+4A^2/\gamma(A),
       \qquad c_A=1+2e^{-A},\\
 J_A(\sigma)&=\sup_{\lambda>0}
  \frac{\log(c_A/c_{A+\lambda\sigma^2})}
       {K(A)+\lambda/\gamma(A)}.
 \end{split}
\]
Choose any $p$ distinct rows and $q$ distinct columns of $D_b$, where $p\ge q\ge1$, and any certified lower singular value $\sigma\ge0$ for that submatrix.
For transfer to the collision rate, require $A>(\log2)e^{2m}$, fix $P,b$ and $\sigma>0$, and first fix $\lambda>0$.
After the input-dependent finite-position onset, take the original word length to infinity, then the singular-value approximation error to zero, and only then the supremum over $\lambda$.
Under these conditions the certificate is the lower-rate value $g(m)+(q/b)J_A(\sigma)$, where
\[
 g(m)=\log\frac{2}{1+2^{1-e^{2m}}}.
\]
Its optimization is over these row and column sets, $q$, $\sigma$, $\lambda$, and $A$ subject to these conditions, within this single limiting $b$-window.
At $\sigma=0$ the increment $J_A(0)$ is zero, so only the baseline $g(m)$ remains.
The rank-energy decomposition gives
\[
 \begin{split}
 D_b&=R_b+E_b,\qquad R_b=C_b-s^2F_b^T,\quad E_b=s^2F_b^T-I_b,\\
 \operatorname{rank}R_b&=\min(b,r),\qquad
 \|E_b\|_F^2\le b(1-e^{-2m}),
 \end{split}
\]
where $\|\cdot\|_F$ is the Frobenius norm.
Here the rank follows from the rank-$r$ finite-Blaschke compression, and the energy bound follows by summing the coefficients of $F$ using $\sum_{j\ge1}c_j^2=s^{-2}-1$.
Restricting to the kernel of the selected part of $R_b$ yields $(q-\min(b,r))_+\sigma^2\le\|E_b\|_F^2$.
The numerator defining $J_A$ is at most $\log c_A$ and at most $2e^{-A}\lambda\sigma^2/c_A$.
Splitting $q$ at $\min(b,r)$ therefore gives, uniformly over the stated optimization,
\[
 \frac qb J_A(\sigma)
 \le\frac{r}{2b}\log2+\frac{1-e^{-2m}}{64}.
\]
If $b\ge2r/m$, the full certificate is at most $g(m)+m\log2/4+m/32$.
Since $g(0)=0$, $g'(0)=\log2$, and $(5/4)\log2+1/32<1$, universal $m_*>0$ and $\epsilon_*>0$ exist such that this value is at most $(1-\epsilon_*)m$ for $0<m<m_*$.
The limiting geometry is taken at fixed $P,b$; its finite-position onset is input dependent.
This is an upper bound on this lower certificate, and gives no upper bound on true entropy or complete pressure.
Shorter scales, joint orbit information and arithmetic constraints on the full neutral law remain distinct questions.
